\akapit Pomiary podzielono na 3 etapy.

\begin{enumerate}
    \item Pomiar parametrów przy wyłączonym tunelu aerodynamicznym i z nominalnym napięciem zasilania silnika
    \item Pomiar parametrów przy włączonym tunelu i z nominalnym napięciem zasilającym
    \item Pomiar parametrów przy wyłączonym tunelu z obniżonym napięciem zasilającym silnik (symulacja rozładowanej baterii)
\end{enumerate}

\akapit W obrębie pomiaru dla jednych warunków dokonywano następujące czynności:

\begin{enumerate}
    \item Zadanie odpowiednich wielkości parametrów:
    \begin{itemize}
        \item  Napięcie zasilania $U_{zaś}$ sterownika zespołu napędowego przy pomocy zasilacza laboratoryjnego
        \item Wysterowanie PWM sterownika zespołu napędowego przy pomocy odpowiedniego zadajnika
        \item Nastawa pracy tunelu aerodynamicznego\footnote{tyko dla drugiego etapu pomiarów}
    \end{itemize}
    \item Odczytanie i zapis wartości mierzonych
    \begin{itemize}
        \item Prąd pobierany $I_{zaś}$ przez zespół napędowy odczytywany z zasilacza laboratoryjnego
        \item Prędkość obrotowa odczytywana z oscyloskopu dzięki układowi czujnika odbiciowego dającego sygnał prostokątny
        \item Siła $T$ tzw. siła ciągu, mierzona przez belkę tensometryczną wraz z mostkiem pomiarowy National Instruments odczytywana za pośrednictwem programu przygotowanego w środowisku LabView
        \item Moment oporowy $M_{opor}$ jako reakcja od momentu obrotowego zespołu napędowego mierzony przez belkę tensometryczną wraz z mostkiem pomiarowy National Instruments odczytywana za pośrednictwem programu przygotowanego w środowisku LabView.
        \item Ciśnienie dynamiczne q strumienia powietrza za śmigłem mierzony przez czujnik HSC Honeywell podłączony do mikrokontrolera integrowanego przez Arduino odczytywane za pośrednictwem programu przygotowanego w środowisku LabView.
        \item Ciśnienie dynamiczne $q_{tun}$ i temperatura t powietrza napływającego w przestrzeni pomiarowej tunelu aerodynamicznego odczytywane ze wskaźnika będącego wyposażeniem tunelu aerodynamicznego.
    \end{itemize}
\end{enumerate}

\akapit Na koniec dokonano obliczeń zgodnie z wzorami przedstawionymi w sekcji \textbf{\ref{sec:podstawy}.Podstawy teoretyczne} oraz przedstawiono otrzymane charakterystyki w sekcji \textbf{\ref{sec:wyniki}.Wyniki}
